\chapter{Funktionenmengen und mengenwertige Funktionen}
\label{b05guide:ch:functionsets}

\section{Funktionen als Elemente einer Menge}

\begin{definition}[Funktionenmenge]
Für Mengen $A,B$ sei
\[
 B^A=\{f\in\mathcal P(A\times B)\mid f\colon A\to B\}.
\]
Die Elemente von $B^A$ sind die Funktionsgraphen mit
Definitionsbereich $A$ und Werten in $B$.
\end{definition}

\begin{theorem}[Existenz und elementare Sonderfälle]
\label{b05guide:thm:functionset}
Für alle Mengen $A,B$ ist $B^A$ eine Menge. Es gelten
\[
 \begin{aligned}
 B^{\varnothing}&=\{\varnothing\},\\
 \varnothing^A&=
 \begin{cases}
  \{\varnothing\},&A=\varnothing,\\
  \varnothing,&A\ne\varnothing,
 \end{cases}\\
 \{b\}^A&=\{A\times\{b\}\},\\
 B^{\{a\}}&=\bigl\{\{(a,b)\}\mid b\in B\bigr\}.
 \end{aligned}
\]
Nach Satz~\ref{b05guide:thm:empty} ist $B^A\ne\varnothing$
genau dann, wenn $A=\varnothing$ oder $B\ne\varnothing$.
\end{theorem}

\begin{theorem}[Zielmengen und Funktionenmengen]
Für beliebige Mengen $A,B,C$ gelten
\[
 \begin{aligned}
 B\subseteq C&\quad\Longrightarrow\quad B^A\subseteq C^A,\\
 B^A\cap C^A&=(B\cap C)^A.
 \end{aligned}
\]
Die Gleichheit wird für die Graphkonvention dieses Leitfadens
verstanden; auch $A=\varnothing$ ist zugelassen.
\end{theorem}

\begin{definition}[Auswertungsfunktion]
Für Mengen $A,B$ ist die Auswertung die Funktion
\[
 \operatorname{ev}_{A,B}\colon B^A\times A\to B,
 \qquad (f,a)\mapsto f(a).
\]
Sie ist für alle $A,B$ definiert. Ist ihr Definitionsbereich leer,
so ist ihr Graph die leere Menge.
\end{definition}

\section{Ein Argument nach dem anderen einsetzen}

\begin{theorem}[Eindeutige Umformung zu funktionswertigen Funktionen]
\label{b05guide:thm:currying}
Seien $A,B,C$ Mengen. Zu jeder Funktion $h\colon A\times B\to C$
existiert genau eine Funktion $H\colon A\to C^B$ mit
\[
 H(a)(b)=h(a,b)\qquad(a\in A,\ b\in B).
\]
Dabei bedeutet $h(a,b)$ den Wert $h((a,b))$, und es gilt
\[
 H(a)=\{(b,h(a,b))\mid b\in B\}\qquad(a\in A).
\]
Umgekehrt bestimmt jede Funktion $H\colon A\to C^B$
genau eine Funktion $h\colon A\times B\to C$ durch dieselbe
Wertgleichung. Beide Konstruktionen gelten auch für leere Mengen.
Die erste wird auch \emph{Currying} genannt.
\end{theorem}

\section{Bild und Urbild als Funktionen auf Potenzmengen}

\begin{definition}[Bildfunktion und Urbildfunktion]
Für $f\colon A\to B$ seien
\[
 \begin{aligned}
 f_*&\colon\mathcal P(A)\to\mathcal P(B),
      &U&\mapsto f[U],\\
 f^*&\colon\mathcal P(B)\to\mathcal P(A),
      &V&\mapsto f^{-1}[V].
 \end{aligned}
\]
Dies sind Funktionen auf Mengen von Teilmengen.
Die Schreibweise $f^*$ bezeichnet die Urbildbildung aus
Kapitel~\ref{b05guide:ch:images}.
\end{definition}

\begin{theorem}[Verkettung der Bild- und Urbildfunktionen]
Für $f\colon A\to B$ und $g\colon B\to C$ gilt
\[
 (g\circ f)_*=g_*\circ f_*,
 \qquad
 (g\circ f)^*=f^*\circ g^*.
\]
Die Verkettung ist im Sinn von
Kapitel~\ref{b05guide:ch:composition} zu lesen.
Die erste Gleichheit gilt für Funktionen
$\mathcal P(A)\to\mathcal P(C)$, die zweite für Funktionen
$\mathcal P(C)\to\mathcal P(A)$.
\end{theorem}

\begin{theorem}[Bestimmtheit durch die Wirkung auf Teilmengen]
\label{b05guide:thm:operator-determination}
Für Funktionen $f,g\colon A\to B$ sind gleichwertig:
\[
 f=g,
 \qquad f_*=g_*,
 \qquad f^*=g^*.
\]
Bereits die Werte auf einelementigen Mengen genügen:
\[
 \begin{aligned}
 f=g
 &\quad\Longleftrightarrow\quad
    \forall a\in A\,\bigl(f[\{a\}]=g[\{a\}]\bigr),\\
 f=g
 &\quad\Longleftrightarrow\quad
    \forall b\in B\,
      \bigl(f^{-1}[\{b\}]=g^{-1}[\{b\}]\bigr).
 \end{aligned}
\]
Die Aussagen gelten auch bei leerem Definitionsbereich.
\end{theorem}

\section{Spuren mengenwertiger Funktionen}

\begin{definition}[Mengenwertige Funktion und Spur]
Eine Funktion $F\colon A\to\mathcal P(B)$ heißt hier
\emph{mengenwertig}; ihr Wert $F(a)$ ist eine Teilmenge von $B$.
Ihre \emph{Spurfunktion} ist
\[
 T_F\colon B\to\mathcal P(A),\qquad
 T_F(b)=\{a\in A\mid b\in F(a)\}.
\]
Damit erfaßt $T_F(b)$ genau diejenigen Argumente, deren
zugeordnete Menge das Element $b$ enthält.
\end{definition}

\begin{theorem}[Rekonstruktion aus den Spuren]
\label{b05guide:thm:traces}
Sei $F\colon A\to\mathcal P(B)$. Für $a\in A$ und $b\in B$ gilt
\[
 b\in F(a)\quad\Longleftrightarrow\quad a\in T_F(b).
\]
Insbesondere gelten
\[
 F(a)=\{b\in B\mid a\in T_F(b)\}\quad(a\in A),
 \qquad T_{T_F}=F.
\]
Bei der zweiten Spurbildung sind $A$ und $B$ vertauscht.
Für $F,G\colon A\to\mathcal P(B)$ gilt außerdem
\[
 F=G\quad\Longleftrightarrow\quad T_F=T_G.
\]
Diese Aussagen umfassen leere Mengen $A,B$ und leere Werte $F(a)$.
\end{theorem}
